% Exact matrix theorem and full reduction to the cited Markov renewal theorem.
\section{Matrix renewal at a locally expanding center}
\label{app:matrix-renewal}

\subsection{Definitions and explicit assumptions}
Let $A_1,A_2,\ldots$ be iid invertible matrices. Work on a compact
communicating class $\mathcal S$ of signed unit directions; projective
space is sufficient when no statement depends on sign. For a nonzero
vector, write $U=r_0e^Z S$. Define the Markov additive process
\begin{equation}
 S_{n+1}=\frac{A_{n+1}S_n}{\|A_{n+1}S_n\|},\qquad
 Z_{n+1}=Z_n+Y_{n+1},\qquad
 Y_{n+1}=\log\|A_{n+1}S_n\|.
 \label{eq:map-exact}
\end{equation}
Its joint direction--increment kernel is $K(s,ds',dy)$. The term
``Markov additive'' means precisely that this kernel depends on the
current direction but not on the current log radius. Let $P_q$ be
\eqref{eq:matrix-pressure-main}, acting on continuous functions with
the supremum norm, and let $k(q)$ be its spectral radius. In particular
$P_0 1=1$ and $k(0)=1$.

\begin{assumption}[Negative-order renewal hypotheses]
\label{ass:matrix-renewal}
The following conditions hold.
\begin{enumerate}
\item For some $\alpha>0$ there exist a continuous $h$, constants
$0<h_{\min}\le h\le h_{\max}<\infty$, and a probability measure $\chi$
such that
\begin{equation}
 P_{-\alpha}h=h,\quad \chi P_{-\alpha}=\chi,\quad
 \chi(h)=1,\quad k(-\alpha)=1.
 \label{eq:matrix-eigen}
\end{equation}
For some $t\in(0,\alpha)$, $P_{-\alpha+t}$ is bounded on
$C(\mathcal S)$ and $k(-\alpha+t)<1$.
\item The probability kernel
\begin{equation}
 \widehat K(s,ds',dy)
   =e^{-\alpha y}\frac{h(s')}{h(s)}K(s,ds',dy)
 \label{eq:matrix-tilt}
\end{equation}
has a positive Harris recurrent angular chain, with invariant law
$\widehat\chi(ds)=h(s)\chi(ds)$, and each one-step angular law
$\widehat P(s,\cdot)$ is absolutely continuous with respect to
$\widehat\chi$. Under $\widehat\chi$,
\begin{equation}
 \widehat{\mathbb E}|Y_1|<\infty,
 \qquad \widehat{\mathbb E}Y_1=-m<0.
 \label{eq:matrix-tilt-drift}
\end{equation}
\item The tilted Markov additive process is nonarithmetic: there is no
$b>0$ and measurable $c:\mathcal S\to[0,b)$ such that
$Y_1\in c(S_0)-c(S_1)+b\mathbb Z$ almost surely under its stationary
angular law.
\end{enumerate}
\end{assumption}

These are hypotheses on the actual random matrices, not conclusions
from a positive one-step scalar log gain. Absolute continuity in the
second condition is a convenient sufficient condition extending the
almost-every-start Markov renewal theorem to every starting direction.
Singular or finitely supported matrix laws may require a different
all-start renewal theorem; the present statement does not silently
cover them. If the dominant eigenvalue and eigenfunction are
differentiable at $-\alpha$, differentiating $P_qh_q=k(q)h_q$ and
integrating against $\chi$ gives
$k'(-\alpha)=\chi(P'_{-\alpha}h)=-m$.
Thus $m=-(\log k)'(-\alpha)$ in that case. Strict convexity of
$\log k$ between the two roots guarantees $k(-\alpha+t)<1$.

Let $\tau=\inf\{n\ge0:Z_n\ge0\}$ and let $\zeta$ be an entry law on
$(-\infty,0)\times\mathcal S$. We impose
\begin{equation}
 L:=\mathbb E_\zeta\tau\in(0,\infty),\qquad
 \mathbb E_\zeta e^{-\alpha Z_0}<\infty.
 \label{eq:matrix-entry}
\end{equation}
At the transition that would exit, the regenerative reference chain
instead draws its next state from an independent copy of $\zeta$.
Consequently each cycle contains states numbered $0,\ldots,\tau-1$.
This construction specifies a reference process. Probabilistic
regeneration of actual SGD is addressed separately below.

\subsection{A uniform occupation bound derived from the spectrum}
\begin{lemma}[Tilted interval potentials are uniformly bounded]
\label{lem:matrix-potential-bound}
Under Assumption~\ref{ass:matrix-renewal}(1), for each $b>0$ there is
$C_b<\infty$ such that
\begin{equation}
 \sup_{z,s,x}\widehat{\mathbb E}_{z,s}
 \sum_{n\ge0}{\bf1}_{\{Z_n\in(x,x+b]\}}\le C_b.
 \label{eq:uniform-renewal-bound}
\end{equation}
\end{lemma}
\begin{proof}
The exact change of measure gives
\begin{equation}
 \widehat{\mathbb E}_s
 e^{t(Z_n-Z_0)}
 =\frac{P_{-\alpha+t}^{,n}h(s)}{h(s)}
 \le C\rho^n
 \label{eq:tilted-exp-contraction}
\end{equation}
for some $\rho<1$ and $C<\infty$.
Indeed, choose $\rho$ strictly between $k(-\alpha+t)$ and $1$;
the spectral-radius formula bounds the operator powers by a constant
times $\rho^n$, and $h$ is bounded above and away from zero.
If $Z_0\in(x,x+b]$, exponential Markov inequality yields
\[
 \widehat{\mathbb P}_{Z_0,s}(Z_n\in(x,x+b])
 \le e^{tb}C\rho^n.
\]
Summing over $n$ bounds occupation by $e^{tb}C/(1-\rho)$.
For any other initial state, apply this bound at its first visit to
the interval, using the strong Markov property. If there is no first
visit the occupation is zero. This proves the claim, uniformly over
both direction and log radius.
\end{proof}

\subsection{Unrestricted and killed renewal asymptotics}
Write $U_{z,s}$ and $\widehat U_{z,s}$ for the unrestricted occupation
measures of \eqref{eq:map-exact} under its original and tilted laws.
For $n$ steps, multiplying \eqref{eq:matrix-tilt} along a path gives
the exact likelihood ratio
\begin{equation}
 \left.\frac{d\widehat{\mathbb P}_{z,s}}
 {d\mathbb P_{z,s}}\right|_{\mathcal F_n}
 = e^{-\alpha(Z_n-z)}\frac{h(S_n)}{h(s)}.
 \label{eq:matrix-rn}
\end{equation}
Consequently
\begin{equation}
 U_{z,s}(dx,ds')=
 e^{\alpha(x-z)}\frac{h(s)}{h(s')}
 \widehat U_{z,s}(dx,ds').
 \label{eq:potential-tilt-exact}
\end{equation}

Apply the Markov renewal theorem of \citet[Theorem 1]{alsmeyer1997}
to the positive-drift process $-Z$. For an angular set $D$ and a finite
log interval, its conclusion, together with the exponential test
weight in \eqref{eq:potential-tilt-exact}, gives
\begin{equation}
 \lim_{x\to-\infty}e^{-\alpha x}U_{z,s}((x,x+b]\times D)
 = e^{-\alpha z}h(s)
 \frac{e^{\alpha b}-1}{\alpha m}\,
 \chi(D).
 \label{eq:unrestricted-matrix-renewal}
\end{equation}
Initially the theorem gives the limit for $\widehat\chi$-almost every
starting direction. To obtain it for an arbitrary $s$, split the
tilted occupation into its time-zero atom and its potential after the
first step. The time-zero atom eventually lies outside the translated
interval. The next direction has a law absolutely continuous with
respect to $\widehat\chi$, so the almost-everywhere limit applies after
this step. Lemma~\ref{lem:matrix-potential-bound}, together with the
bounded test weight $e^{\alpha u}/h(s')$, permits dominated convergence
over this first step. This proves the limit for every $s$ and every
finite initial $z$. The angular factor is $\chi$, since division by
$h(s')$ cancels the $h(s')$ in $\widehat\chi$.

Moreover, \eqref{eq:potential-tilt-exact} and the same lemma give a
constant $C_b'$ such that, for all $x,z,s$,
\begin{equation}
 e^{-\alpha x}U_{z,s}((x,x+b]\times\mathcal S)
 \le C_b'e^{-\alpha z}h(s).
 \label{eq:original-renewal-domination}
\end{equation}
This bound is the domination needed for integration over entry and exit
states. It has been proved, rather than included as an unverified
``uniform renewal'' assumption.

\begin{proof}[Proof of Theorem~\ref{thm:matrix-central-main}]
Let $G_\zeta$ denote expected occupation before $\tau$.
The strong Markov property gives, on bounded log intervals,
\begin{equation}
 G_\zeta=U_\zeta-
 \mathbb E_\zeta U_{Z_\tau,S_\tau}.
 \label{eq:killed-subtraction}
\end{equation}
Finite mean exit time implies $\tau<\infty$ almost surely under
$\mathbb P_\zeta$. The entry moment in \eqref{eq:matrix-entry}, bounded
$h$, and $Z_\tau\ge0$ permit integrating
\eqref{eq:unrestricted-matrix-renewal} in both terms by
\eqref{eq:original-renewal-domination}. Hence
\begin{equation}
 \lim_{x\to-\infty}e^{-\alpha x}G_\zeta((x,x+b]\times D)
 = D_\alpha\frac{e^{\alpha b}-1}{\alpha m}\chi(D),
 \label{eq:killed-asymptotic}
\end{equation}
where
\begin{equation}
 D_\alpha=
 \mathbb E_\zeta\!\left[e^{-\alpha Z_0}h(S_0)
                 -e^{-\alpha Z_\tau}h(S_\tau)\right].
 \label{eq:matrix-central-constant}
\end{equation}
Its sign cannot be read by comparing endpoint radii alone, because
directions change. For each fixed start, summing the stopped likelihood
ratio \eqref{eq:matrix-rn} over events $\{\tau=n\}$ instead gives
\[
 \mathbb E_{z,s}[e^{-\alpha Z_\tau}h(S_\tau);\tau<\infty]
 =e^{-\alpha z}h(s)\,
   \widehat{\mathbb P}_{z,s}(\tau<\infty).
\]
Therefore
\begin{equation}
 D_\alpha=
 \mathbb E_\zeta\!\left[
 e^{-\alpha Z_0}h(S_0)
 \widehat{\mathbb P}_{Z_0,S_0}(\tau=\infty)\right].
 \label{eq:positive-constant}
\end{equation}
The stated positive-survival condition makes $D_\alpha>0$.
It is enough for a positive entry mass to be sufficiently far below
the upper boundary: by \eqref{eq:tilted-exp-contraction},
\[
 \widehat{\mathbb P}_{z,s}(\tau<\infty)
 \le \sum_{n\ge1}\widehat{\mathbb P}_{z,s}(Z_n\ge0)
 \le \frac{C\rho}{1-\rho}e^{tz},
\]
which is strictly below one for all sufficiently negative $z$.

The regenerative occupation formula
\begin{equation}
 \pi(F)=\frac{1}{L}\,
 \mathbb E_\zeta\sum_{n=0}^{\tau-1}{\bf1}_{\{(Z_n,S_n)\in F\}}
 \label{eq:regeneration-formula}
\end{equation}
defines an invariant probability. For a bounded test function, summing
its expected one-step change over a cycle telescopes to the difference
between two cycle starts, both with law $\zeta$; hence its integral
against $\pi P-\pi$ is zero.
To pass from intervals to the half-line $(-\infty,x]$, partition it
into intervals $(x-(j+1)b,x-jb]$. The bound
\eqref{eq:original-renewal-domination} dominates their rescaled
occupations by a summable geometric sequence in $j$. Summing
\eqref{eq:killed-asymptotic} therefore yields
\begin{equation}
 \lim_{x\to-\infty}e^{-\alpha x}\pi(Z\le x,S\in D)
 = \frac{D_\alpha}{\alpha m L}\chi(D).
 \label{eq:matrix-final-constant}
\end{equation}
Substituting $x=\log(r/r_0)$ proves the theorem, with
$C=D_\alpha/(\alpha mL)$.
\end{proof}

The theorem proves a small-ball asymptotic. Differentiating an
asymptotic cumulative distribution without further hypotheses is
invalid. A radial-density conclusion requires the corresponding
renewal-density result, for example a suitable smoothing condition
giving convergence of local densities. Angular density conclusions
additionally require regularity of $\chi$.

\subsection{A noncommuting law that verifies every renewal hypothesis}
\label{app:realizable-matrix-renewal}
The preceding assumptions are realizable even by symmetric random
matrices. For $d\ge2$, take
\begin{equation}
 A=M(I-2vv^\top),\qquad
 v\sim\operatorname{Unif}(S^{d-1}),\qquad
 M\sim\tfrac14\operatorname{Unif}[.49,.51]
      +\tfrac34\operatorname{Unif}[1.99,2.01],
 \label{eq:reflection-example}
\end{equation}
with independent draws at each update. Independent reflections
generically do not commute. Their norm gain is $M$, so
$P_q=m(q)R$, where $m(q)=\mathbb EM^q$ and $R$ is the angular
reflection kernel. Thus $h=1$, $k(q)=m(q)$, and $\chi$ is uniform
surface probability. This scalar formula is special to this example;
the matrix theorem itself retains the coupled angular dynamics.

For completeness the reflection kernel has, with respect to uniform
surface probability, density
\begin{equation}
 R(s,ds')=\|s-s'\|^{2-d}\chi(ds')\quad(s'\ne s).
 \label{eq:reflection-angular-density}
\end{equation}
To verify it, put $s=e_1$ and let $\theta$ be the angle between $s$ and
$s'$. Reflection gives $|v^\top s|=\sin(\theta/2)$.
Changing variables from the spherical coordinate of $v$ gives
the angle density $c_d\cos^{d-2}(\theta/2)$, whereas a uniform
$s'$ has angle density $c_d\sin^{d-2}\theta$. Their ratio is
$(2\sin(\theta/2))^{2-d}=\|s-s'\|^{2-d}$.
In particular $R(s,\cdot)\ge2^{2-d}\chi(\cdot)$, so it is a
uniformly mixing Harris kernel, and its symmetry verifies invariance
of $\chi$. Its one-step laws are absolutely continuous.

The top product norm equals $\prod_{j=1}^n M_j$, since every remaining
factor is orthogonal. Hence
\[
 \gamma=\mathbb E\log M
 \ge\tfrac14\log(.49)+\tfrac34\log(1.99)>0,
 \qquad
 m(-1)\le\frac{1}{4(.49)}+\frac{3}{4(1.99)}<1.
\]
Strict log convexity and the positive mass below one imply a unique
negative root $m(-\alpha)=1$ with $\alpha>1$ and
$m(-\alpha+t)<1$ for every $t\in(0,\alpha)$.
Tilting changes only the scalar multiplier law; it leaves $R$
unchanged. The tilted log increment has bounded support, a density,
and strictly negative mean $m'(-\alpha)<0$, so all spectral and
nonarithmetic hypotheses hold.

Choose the entry log radius $Z_0=-1$ and an independent uniform
direction. Since the original log increments are iid with mean
$\gamma>0$ and are bounded above by $\log2.01$, Wald's identity at
$\tau\wedge n$ gives
\[
 \gamma\mathbb E(\tau\wedge n)
   =\mathbb EZ_{\tau\wedge n}+1
   \le\log2.01+1.
\]
The strong law gives almost-sure exit; monotone convergence now gives
finite mean exit time. The entry inverse moment is $e^\alpha$.
Under the tilted law $e^{\alpha Z_n}$ is a nonnegative martingale;
optional stopping yields
$\widehat{\mathbb P}_{-1,s}(\tau<\infty)\le e^{-\alpha}<1$.
The central-law coefficient is therefore strictly positive.
All hypotheses of Theorem~\ref{thm:matrix-central-main} have been
verified for this noncommuting law.

The matrices are symmetric, so they can be written
$A=I-\eta H$ for symmetric sample Hessians. In $d=2$,
$\mathbb E(I-2vv^\top)=0$, so $\mathbb EH=I/\eta$ is positive
definite even though individual sample Hessians need not be. The
independent restart remains an explicit return mechanism of the
reference process; it is not silently identified with an SGD update.

\subsection{An exact, conditional transfer statement for nonlinear SGD}
\begin{proposition}[Nonlinear occupation transfer]
\label{prop:nonlinear-transfer}
Consider a nonlinear chain that has genuine regenerative cycles of
finite mean length $T$, with cycle endpoints outside the local domain
$\mathcal D=\{(z,s):z<0\}$. Let $Q$ be its actual transition kernel
killed on leaving $\mathcal D$, and let $L_0$ be the kernel of
\eqref{eq:map-exact} killed at its upper exit. Let $G$ be the expected
measure of all local visits in one complete cycle, and $\nu$ the
expected measure of every entry into $\mathcal D$ in that cycle. Put
\begin{equation}
 \mathcal J=\nu+G(Q-L_0),\qquad
 H(z,s)=e^{-\alpha z}h(s)
       \widehat{\mathbb P}_{z,s}(\tau=\infty).
 \label{eq:nonlinear-source}
\end{equation}
Assume the killed linear chain exits from $G$-almost every starting
point and that
\begin{equation}
 \int e^{-\alpha z}\,|\mathcal J|(dz,ds)<\infty,
 \qquad D_Q:=\int H\,d\mathcal J>0.
 \label{eq:source-integrability}
\end{equation}
Under Assumption~\ref{ass:matrix-renewal}, the genuine invariant law
has the same central exponent:
\begin{equation}
 \lim_{r\downarrow0}
 \frac{\pi(\|U\|\le r,S\in D)}{(r/r_0)^\alpha}
 = \frac{D_Q}{\alpha m\mathbb ET}\chi(D).
 \label{eq:nonlinear-central-law}
\end{equation}
\end{proposition}
\begin{proof}
Counting each visit as either an entry or the successor of a previous
local visit yields $G=\nu+GQ$. Thus $G=\mathcal J+GL_0$ exactly, and
\[
 \sum_{j=0}^{n-1}\mathcal J L_0^j=G-GL_0^n.
\]
Since $G$ is finite and the killed linear chain exits $G$-almost
everywhere, $GL_0^n$ tends to zero in total variation. For each fixed
bounded log interval, the killed potential is bounded by its
unrestricted potential, and \eqref{eq:original-renewal-domination}
shows absolute integrability against $|\mathcal J|$ by
\eqref{eq:source-integrability}. We may therefore pass to the infinite
sum and then integrate the killed asymptotic against the signed source.
The coefficient is $D_Q$. Summing intervals as in the preceding proof
and dividing occupation by the complete nonlinear cycle length proves
\eqref{eq:nonlinear-central-law}.
\end{proof}

The source is signed. A nonnegative entry law does not guarantee
$D_Q>0$, because nonlinear correction can cancel its leading term.
The source condition can be checked without assuming the desired
central exponent. For example, suppose $V_p(z,s)$ is comparable to
$e^{-pz}$ for some $0<p<\alpha$ and satisfies
\begin{equation}
 QV_p\le\rho_p V_p,\qquad \rho_p<1,
 \qquad \nu(e^{-\alpha z})<\infty,
 \label{eq:inverse-occupation-drift}
\end{equation}
and
\begin{equation}
 \int e^{-\alpha z'}
 |Q-L_0|((z,s),d(z',s'))\le C V_p(z,s).
 \label{eq:weighted-kernel-discrepancy}
\end{equation}
Decomposing the cycle into its actual local excursions gives
$G=\nu\sum_{n\ge0}Q^n$ and hence
$G(V_p)\le\nu(V_p)/(1-\rho_p)<\infty$.
Equations \eqref{eq:inverse-occupation-drift}--\eqref{eq:weighted-kernel-discrepancy}
then imply \eqref{eq:source-integrability}'s integrability assertion.
An eigenfunction $h_p$ with $P_{-p}h_p=k(-p)h_p$ and $k(-p)<1$
suggests $V_p=e^{-pz}h_p(s)$ for this independent drift calculation.

A Taylor expansion $F_B(u)=A_Bu+O(\|u\|^{1+\delta})$ does not by
itself prove \eqref{eq:weighted-kernel-discrepancy}: nearby point masses
have order-one total variation distance. Smooth transition laws or a
different perturbative renewal theorem are needed. Fixed nondegenerate
additive noise also destroys the multiplicative approximation as
$r\downarrow0$. For that case an intermediate small-noise annulus,
with estimates uniform in the noise scale, must replace the literal
zero-radius limit. We do not claim these transfer conditions for the
neural experiments.
